\title{A Pattern Collection for Generating Accessible Teaching Materials  in \ac{stem}}

\titlerunning{Accessible Teaching Materials in \ac{stem}}

\author{Diethelm Bienhaus\inst{1}\orcidID{0000-0002-9207-3412}  
\and \\
Michael Kreutzer\inst{1}\orcidID{0000-0003-0748-7707}
\and \\
Florian von Zabiensky\inst{1}\orcidID{0000-0003-0951-0288}}
%
\authorrunning{D. Bienhaus, M. Kreutzer, F. v. Zabiensky}



\institute{Technische Hochschule Mittelhessen University of Applied Sciences, Germany
\email{\{diethelm.bienhaus, michael.kreutzer\}@mni.thm.de},
\\
Website: \url{mni.thm.de}
}
%
\maketitle              % typeset the header of the contribution
%

\begin{abstract}
\ac{stem} education relies heavily on visual artifacts such as diagrams, circuit schematics, plots, tables, formulas, and graphical simulation tools. These artifacts often remain difficult or impossible to access for \ac{bvi} students when they are only provided as visual representations. At the same time, long and densely structured academic materials can impose additional barriers for learners with impairments such as \ac{adhd} or dyslexia, who may benefit from segmented and multimodal access.

This paper extends the EuroPLoP 2025 pattern collection \textit{A Pattern Collection for Generating Accessible Teaching Materials for Blind and Visually Impaired Students in Computer Science and Electrical Engineering}. The earlier collection focused on transforming visual \ac{stem} artifacts into accessible textual, tactile, or auditory representations. We add two AI-mediated patterns that address accessibility during actual learning situations rather than only during material preparation.

\textsc{Accessibility-Aware AI Tutor} describes how to provide persistent, course-specific, and accessibility-aware support that is grounded in curated course materials, explicit accessibility rules, and repeatable transformation workflows. \textsc{Structured Audio Companion} describes how to provide segmented, source-grounded audio orientation with transcripts, timecodes, and explicit pointers to precise accessible representations of visual artifacts.

Together, the two patterns show that generative artificial intelligence can support accessibility when it is embedded in course context, stable workflows, and human-in-the-loop quality assurance. AI-generated explanations and audio summaries should not replace precise accessible representations of visual \ac{stem} artifacts; rather, they should help learners find, understand, and use them.

\keywords{Accessibility \and Blind and Visually Impaired Learners \and STEM Education \and Generative Artificial Intelligence \and Large Language Models \and ADHD \and Dyslexia}
\end{abstract}

% Keywords: Accessibility, Blind and Visually impaired Learners, Large Language Models,  Attention-Deficit/Hyperactivity Disorder (ADHD), Dyslexia